\par\Needspace{4\baselineskip}
\begin{center}
\SourceLeaf{027c}\pdfbookmark[2]{10.}{article-010}10.
\end{center}
The equations in Article 8 contain everything needed to compute the eccentric and mean
anomalies from the true anomaly, or the eccentric and true anomalies from the mean
anomaly. Formula VII is commonly used to derive the eccentric anomaly from the true
anomaly. For this purpose, however, equation X is usually preferable, especially when
the eccentricity is not too large, in which case $E$ can be computed more precisely
by X than by VII. Moreover, when equation X is used, the logarithm of the sine of $E$
required in XII is obtained directly from equation VIII. If VII were used, this
logarithm would have to be taken from the tables. Thus, if the logarithm is also
taken from the tables in the former method, the comparison confirms that the calculation
was carried out correctly. Such checks and confirmations of calculations are always
of great value. We shall therefore take constant care to provide for them in all
the methods presented in this work wherever this can conveniently be done.
--- For further illustration, we give a fully worked example.

\Needspace{22\baselineskip}
Let $v = 310^0\,55^{\prime}\,29^{\prime\prime}{,}64$, $\varphi = 14^0\,12^{\prime}\,1^{\prime\prime}{,}87$, and $\log r = 0{,}3307640$ be given. We seek $p$, $a$, $E$ and $M$.

\SourceLeaf{028}
\begin{center}
\begin{tabular}{@{}ll@{}}
\makebox[16em][l]{$\log\sin\varphi$\hfill .\,.\,.\,.\,.\,.\,.\,.\,.\,.} & $9{,}389\,7262$\\[.3em]
\makebox[16em][l]{$\log\cos v$\hfill .\,.\,.\,.\,.\,.\,.\,.\,.\,.} & $9{,}816\,2872$\\[.3em]
\hline
 & $9{,}206\,0134$\\[.3em]
\multicolumn{2}{r}{whence $e\cos v=0{,}1606991$}\\[.3em]
\hline
\makebox[16em][l]{$\log(1+e\cos v)$\hfill .\,.\,.\,.\,.\,.} & $0{,}064\,7197$\\[.3em]
\makebox[16em][l]{$\log r$\hfill .\,.\,.\,.\,.\,.\,.\,.\,.\,.\,.\,.\,.} & $0{,}330\,7640$\\[.3em]
\hline
\makebox[16em][l]{$\log p$\hfill .\,.\,.\,.\,.\,.\,.\,.\,.\,.\,.\,.\,.} & $0{,}395\,4837$\\[.3em]
\makebox[16em][l]{$\log\cos\varphi^2$\hfill .\,.\,.\,.\,.\,.\,.\,.\,.} & $9{,}973\,0448$\\[.3em]
\hline
\makebox[16em][l]{$\log a$\hfill .\,.\,.\,.\,.\,.\,.\,.\,.\,.\,.\,.\,.} & $0{,}422\,4389$\\[.3em]
\hline
\makebox[16em][l]{$\log\sin v$\hfill .\,.\,.\,.\,.\,.\,.\,.\,.\,.} & $9{,}878\,2740_n$\footnotemark\\[.3em]
\makebox[16em][l]{$\log\BareRoot{\frac{p}{r}}$\hfill .\,.\,.\,.\,.\,.\,.\,.\,.\,.} & $0{,}032\,3598.5$\\[.3em]
\hline
 & $9{,}845\,9141.5_n$\\[.3em]
\makebox[16em][l]{$\log\sin\frac{1}{2}\varphi$\hfill .\,.\,.\,.\,.\,.\,.\,.\,.} & $9{,}092\,0395$\\[.3em]
\hline
\noalign{\vskip2pt}
\makebox[16em][l]{$\log\sin\frac{1}{2}(v-E)$\hfill .\,.\,.\,.\,.} & $8{,}937\,9536.5_n$\\[.3em]
\end{tabular}
\end{center}
\footnotetext{The letter $n$ appended to a logarithm indicates that the corresponding number is negative.}

\noindent Hence
\[
\begin{gathered}
\frac{1}{2}(v-E)=-4^0\,58^{\prime}\,22^{\prime\prime}{,}94;\quad v-E=-9^0\,56^{\prime}\,45^{\prime\prime}{,}88;\\
E=320^0\,52^{\prime}\,15^{\prime\prime}{,}52.
\end{gathered}
\]
\noindent We also obtain
\begin{center}
\begin{tabular}{@{}ll@{}}
\makebox[16em][l]{$\log e$\hfill .\,.\,.\,.\,.\,.\,.\,.\,.} & $9{,}389\,7262$\\[.3em]
\makebox[16em][l]{$\log206264{,}8$\hfill .\,.\,.} & $5{,}314\,4251$\\[.3em]
\hline
\makebox[16em][l]{$\log e$ in arcsec:\hfill .\,.\,.} & $4{,}704\,1513$\\[.3em]
\makebox[16em][l]{$\log\sin E$\hfill .\,.\,.\,.\,.\,.} & $9{,}800\,0767_n$\\[.3em]
\hline
 & $4{,}504\,2280_n$\\[.3em]
\end{tabular}
\end{center}

\begin{center}
Calculation of $\log\sin E$ by formula VIII.
\end{center}
\begin{center}
\begin{tabular}{@{}ll@{}}
\makebox[16em][l]{$\log\frac{r}{p}\sin v$\hfill .\,.\,.\,.\,.} & $9{,}813\,5543_n$\\[.3em]
\makebox[16em][l]{$\log\cos\varphi$\hfill .\,.\,.\,.\,.\,.} & $9{,}986\,5224$\\[.3em]
\hline
\makebox[16em][l]{$\log\sin E$\hfill .\,.\,.\,.\,.\,.} & $9{,}800\,0767_n$\\[.3em]
\hline
\end{tabular}
\end{center}

\noindent Hence
\[
\begin{gathered}
e\sin E\ \mathrm{in\ arcseconds}=-31932^{\prime\prime}{,}14=-8^0\,52^{\prime}\,12^{\prime\prime}{,}14\\
\mathrm{and}\quad M=329^0\,44^{\prime}\,27^{\prime\prime}{,}66.
\end{gathered}
\]
\Needspace{12\baselineskip}
\noindent Using formula VII, the calculation of $E$ would proceed as follows.
\begin{center}
\begin{tabular}{@{}ll@{}}
\multicolumn{2}{l}{$\frac{1}{2}v=155^0\,27^{\prime}\,44^{\prime\prime}{,}82$}\\[.3em]
\makebox[16em][l]{$\log\operatorname{tang}\frac{1}{2}v$\hfill .\,.\,.\,.\,.\,.\,.\,.} & $9{,}659\,4579_n$\\[.3em]
\multicolumn{2}{l}{$45^0-\frac{1}{2}\varphi=37^0\,53^{\prime}\,59^{\prime\prime}{,}065$}\\[.3em]
\makebox[16em][l]{$\log\operatorname{tang}(45^0-\frac{1}{2}\varphi)$\hfill .\,.\,.} & $9{,}891\,2427$\\[.3em]
\hline
\noalign{\vskip2pt}
\makebox[16em][l]{$\log\operatorname{tang}\frac{1}{2}E$\hfill .\,.\,.\,.\,.\,.\,.} & $9{,}550\,7006_n$\\[.3em]
\end{tabular}
\end{center}
This gives $\frac{1}{2}E=160^0\,26^{\prime}\,7^{\prime\prime}{,}76$ and $E=320^0\,52^{\prime}\,15^{\prime\prime}{,}52$, as above.
